\chapter{Annexes}
\label{ch:annexes}

\section{Annex A: Full Bill of Materials}
\label{annex:bom}

The full BOM with supplier links is provided in the accompanying spreadsheet file (\texttt{bom\_ssms.xlsx}). A condensed version was presented in Table~\ref{tab:bom}.

\section{Annex B: Complete MicroPython Source Code}
\label{annex:fw-code}

The complete \texttt{main.py} firmware file is included below:

\begin{lstlisting}[language=Python, caption=Complete MicroPython main.py, label=lst:fw-main]
import network
import time
from machine import Pin
from neopixel import NeoPixel
from umqtt.simple import MQTTClient
import json

# Configuration
LED_COUNT = 216
LEDS_PER_BIN = 6
NUM_BINS = 36
MQTT_BROKER = "192.168.1.100"
MQTT_TOPIC_CMD = b"ssms/command"
MQTT_TOPIC_STATUS = b"ssms/status"
WIFI_SSID = "SSMS_NET"
WIFI_PASS = "ssms2025"

# NeoPixel initialisation
data_pin = Pin(18, Pin.OUT)
np = NeoPixel(data_pin, LED_COUNT)

# State variables
current_state = "IDLE"
blink_active = False
blink_bin = 0
blink_color = (0, 0, 255)
blink_period = 500
last_toggle = 0
complete_start = 0

def set_bin(bin_idx, color):
    base = bin_idx * LEDS_PER_BIN
    for i in range(LEDS_PER_BIN):
        np[base + i] = color
    np.write()

def clear_bin(bin_idx):
    set_bin(bin_idx, (0, 0, 0))

def clear_all():
    for i in range(LED_COUNT):
        np[i] = (0, 0, 0)
    np.write()

def mqtt_callback(topic, msg):
    global current_state, blink_active, blink_bin, blink_color, blink_period, complete_start
    payload = msg.decode()
    parts = payload.split(":")
    cmd = parts[0]

    if cmd == "LED_ON":
        idx = int(parts[1])
        r, g, b = int(parts[2]), int(parts[3]), int(parts[4])
        set_bin(idx, (r, g, b))
        current_state = "ACTIVE"
        client.publish(MQTT_TOPIC_STATUS, f"OK:LED_ON:{idx}")
    elif cmd == "LED_OFF":
        idx = int(parts[1])
        clear_bin(idx)
        current_state = "IDLE"
        client.publish(MQTT_TOPIC_STATUS, f"OK:LED_OFF:{idx}")
    elif cmd == "ALL_OFF":
        clear_all()
        current_state = "IDLE"
        blink_active = False
        client.publish(MQTT_TOPIC_STATUS, "OK:ALL_OFF")
    elif cmd == "ALL_ON":
        r, g, b = int(parts[1]), int(parts[2]), int(parts[3])
        for i in range(LED_COUNT):
            np[i] = (r, g, b)
        np.write()
        client.publish(MQTT_TOPIC_STATUS, "OK:ALL_ON")
    elif cmd == "BLINK":
        blink_bin = int(parts[1])
        blink_color = (int(parts[2]), int(parts[3]), int(parts[4]))
        blink_period = int(parts[5])
        blink_active = True
        last_toggle = time.ticks_ms()
        current_state = "SCANNING"
        set_bin(blink_bin, blink_color)
        client.publish(MQTT_TOPIC_STATUS, f"OK:BLINK:{blink_bin}")

def connect_wifi():
    wlan = network.WLAN(network.STA_IF)
    wlan.active(True)
    wlan.connect(WIFI_SSID, WIFI_PASS)
    while not wlan.isconnected():
        time.sleep(0.5)
    print("WiFi connected:", wlan.ifconfig())

def connect_mqtt():
    global client
    client = MQTTClient("ssms_esp32", MQTT_BROKER)
    client.set_callback(mqtt_callback)
    client.connect()
    client.subscribe(MQTT_TOPIC_CMD)
    print("MQTT connected")
    client.publish(MQTT_TOPIC_STATUS, "ONLINE")

connect_wifi()
connect_mqtt()
clear_all()

while True:
    try:
        client.check_msg()

        now = time.ticks_ms()
        if blink_active and current_state == "SCANNING":
            if now - last_toggle >= blink_period // 2:
                last_toggle = now
                base = blink_bin * LEDS_PER_BIN
                if np[base] == (0, 0, 0):
                    for i in range(LEDS_PER_BIN):
                        np[base + i] = blink_color
                else:
                    for i in range(LEDS_PER_BIN):
                        np[base + i] = (0, 0, 0)
                np.write()
        elif current_state == "COMPLETE":
            if now - complete_start >= 2000:
                clear_bin(blink_bin)
                current_state = "IDLE"

        time.sleep_ms(10)
    except OSError:
        time.sleep(5)
        connect_mqtt()
\end{lstlisting}

\section{Annex C: Complete Desktop Application Source Code}
\label{annex:app-code}

The complete desktop application source code is too extensive to reproduce in full here. The file is provided as \texttt{ui\_main.py} in the project archive. Key excerpts are presented in Chapter~\ref{ch:implementation}.

\section{Annex D: MQTT Topic and Command Reference}
\label{annex:mqtt-ref}

\begin{longtable}{p{3cm} p{3cm} p{6cm}}
  \caption{MQTT command reference.}\label{tab:mqtt-ref}\\
  \toprule
  \textbf{Topic} & \textbf{Command} & \textbf{Description} \\
  \midrule
  \endfirsthead
  \multicolumn{3}{c}{\tablename\ \thetable\ --- Continued} \\
  \toprule
  \textbf{Topic} & \textbf{Command} & \textbf{Description} \\
  \midrule
  \endhead
  \bottomrule
  \endfoot
  \texttt{ssms/command} & \texttt{LED\_ON:idx:r:g:b} & Set bin at index to colour \\
  \texttt{ssms/command} & \texttt{LED\_OFF:idx} & Turn off bin at index \\
  \texttt{ssms/command} & \texttt{ALL\_ON:r:g:b} & Set all LEDs to colour \\
  \texttt{ssms/command} & \texttt{ALL\_OFF} & Turn off all LEDs \\
  \texttt{ssms/command} & \texttt{BLINK:idx:r:g:b:ms} & Blink bin at period \\
  \texttt{ssms/status} & \texttt{OK:...} & Acknowledgment \\
  \texttt{ssms/status} & \texttt{ONLINE} & ESP32 boot notification \\
  \texttt{ssms/status} & \texttt{ERROR:msg} & Error message \\
\end{longtable}

\section{Annex E: KiCad Schematic}
\label{annex:schematic}

The full KiCad schematic is provided as a separate PDF file (\texttt{schematic\_ssms.pdf}).

\textit{[Placeholder: Insert full-page KiCad schematic figure here]}

\section{Annex F: LTspice Netlist}
\label{annex:ltspice}

\begin{lstlisting}[caption=LTspice netlist for level-shifter simulation, label=lst:ltspice]
* Level Shifter Simulation: ESP32 GPIO -> 74HCT245 -> WS2813
V1 GPIO_OUT 0 PULSE(0 3.3 0 5n 5n 500n 1.25u)
XU1 A1 B1 VCC GND 74HCT245
RB1 DATA_LINE 0 100
C1 DATA_LINE 0 10p
VCC VCC 0 5
.tran 10u
.end
\end{lstlisting}

\section{Annex G: Acronyms and Abbreviations}
\label{annex:acronyms}

\begin{longtable}{p{3cm} p{9cm}}
  \caption{Acronyms and abbreviations.}\label{tab:acronyms}\\
  \toprule
  \textbf{Acronym} & \textbf{Full Form} \\
  \midrule
  \endfirsthead
  \multicolumn{2}{c}{\tablename\ \thetable\ --- Continued} \\
  \toprule
  \textbf{Acronym} & \textbf{Full Form} \\
  \midrule
  \endhead
  \bottomrule
  \endfoot
  BOM & Bill of Materials \\
  CSV & Comma-Separated Values \\
  EDA & Electronic Design Automation \\
  ERP & Enterprise Resource Planning \\
  ESP32 & Espressif Systems 32-bit MCU \\
  GPIO & General-Purpose Input/Output \\
  GUI & Graphical User Interface \\
  I2C & Inter-Integrated Circuit \\
  LAN & Local Area Network \\
  LED & Light-Emitting Diode \\
  MCU & Microcontroller Unit \\
  MQTT & Message Queuing Telemetry Transport \\
  MVC & Model-View-Controller \\
  PCB & Printed Circuit Board \\
  PFE & Projet de Fin d'\'Etudes \\
  PSU & Power Supply Unit \\
  PUT & Put-away (stocking) operation \\
  QoS & Quality of Service \\
  RFID & Radio-Frequency Identification \\
  RTM & Requirements Traceability Matrix \\
  SSMS & Smart Storage Management System \\
  TVS & Transient Voltage Suppression \\
  UART & Universal Asynchronous Receiver-Transmitter \\
  WMS & Warehouse Management System \\
\end{longtable}

\section{Annex H: References}
\label{annex:references}

The complete list of references is provided in the bibliography section at the end of this document. Key references include:

\begin{enumerate}
  \item Hercog et al., ``Pick-to-Light Systems: A Review and Classification,'' \textit{Sensors}, vol. 22, no. 24, p. 9769, 2022. DOI: 10.3390/s22249769.
  \item Espressif Systems, \textit{ESP32 Technical Reference Manual}, 2023.
  \item Worldsemi, \textit{WS2813 Addressable RGB LED Datasheet}, 2019.
  \item Texas Instruments, \textit{SN74HCT573 Octal Transparent D-Type Latch Datasheet}, 2022.
  \item Hi-Link, \textit{HLK-PM01 AC-DC Power Module Datasheet}, 2020.
  \item OASIS, \textit{MQTT Version 3.1.1}, 2014.
  \item MicroPython, \textit{MicroPython Documentation}, 2023.
  \item PyInstaller, \textit{PyInstaller Documentation}, 2023.
  \item KiCad, \textit{KiCad EDA Documentation}, 2023.
  \item G\"unter et al., ``Industry 4.0 in Warehousing: A Systematic Literature Review,'' \textit{Int. Journal of Production Research}, vol. 58, no. 10, pp. 2984--3010, 2020.
\end{enumerate}
